\section*{अध्याय \ref{ch:g11:ocean-ridges} --- महासागरीय कटक और महासागरीय स्थलमंडल}

\begin{solution}{exo:g11:ocean-ridges:1}
\omterm{def:g11:ocean-ridges:peridotite}{पेरिडोटाइट} एक गहरे रंग की, मोटे कणों वाली चट्टान है जो मुख्यतः ओलिवीन से बनी है, कुछ
पाइरॉक्सीन के साथ; प्रावार का अधिकांश भाग इससे बना है। \omterm{def:g11:ocean-ridges:gabbro}{गैब्रो} प्लेजियोक्लेज़ फ़ेल्डस्पार और पाइरॉक्सीन की
एक गहरे रंग की, मोटे कणों वाली आग्नेय चट्टान है, जो गहरी \omterm{def:g11:ocean-ridges:oceanic-lithosphere}{महासागरीय भूपर्पटी} की सबसे प्रचुर
चट्टान है। \omterm{def:g11:ocean-ridges:gabbro}{गैब्रो} भूपर्पटी का अंग है, \omterm{def:g11:ocean-ridges:peridotite}{पेरिडोटाइट} प्रावार का।
\omterm{def:g11:ocean-ridges:gabbro}{गैब्रो} समुद्री तल के बहुत नीचे धीरे-धीरे क्रिस्टलित हुआ, जहाँ मैग्मा
धीरे ठंडा हुआ: उसके क्रिस्टलों को बढ़ने का समय मिला। बेसाल्ट लावा समुद्री
तल पर जल्दी ठंडा हुआ, और उसके क्रिस्टल छोटे रह गए।
\end{solution}

\begin{solution}{exo:g11:ocean-ridges:2}
समुद्री तल से नीचे की ओर: अवसाद (अक्ष पर कोई नहीं), बेसाल्ट लावा (तकिए
और चादरी प्रवाह), चादरी डाइक, \omterm{def:g11:ocean-ridges:gabbro}{गैब्रो}, फिर \omterm{def:g11:seismic-structure:moho}{मोहो} के नीचे
प्रावार का \omterm{def:g11:ocean-ridges:peridotite}{पेरिडोटाइट}। छेद ने $1507/7100 \approx 0.21$ पार किया, भूपर्पटी का
लगभग पाँचवाँ भाग: गहरे \omterm{def:g11:ocean-ridges:gabbro}{गैब्रो} और \omterm{def:g11:seismic-structure:moho}{मोहो} तक कोई वेधन कभी नहीं पहुँचा।
\end{solution}

\begin{solution}{exo:g11:ocean-ridges:3}
$d(0) = \qty{2500}{m}$; $d(9) = 2500 + 350 \times 3 = \qty{3550}{m}$;
$d(25) = 2500 + 350 \times 5 = \qty{4250}{m}$।
\end{solution}

\begin{solution}{exo:g11:ocean-ridges:4}
\omterm{def:g11:ocean-ridges:normal-fault}{सामान्य भ्रंश} एक झुका हुआ भ्रंश है जिसके अनुदिश भ्रंश-तल के ऊपर का
खंड नीचे के खंड के सापेक्ष नीचे खिसका है। धीमे कटक पर प्लेटें
ठंडे, भंगुर \omterm{def:g11:lithosphere-asthenosphere:lithosphere}{स्थलमंडल} को खींचकर अलग करती हैं: वह खिंचता है, और टूट जाता है।
भ्रंशों के ऊपर के खंड नीचे, अक्ष की ओर, सरकते हैं, और अक्षीय घाटी
वह खंड है जो उनके बीच धँस गया है।
\end{solution}

\begin{solution}{exo:g11:ocean-ridges:5}
यह दाब के कारण नहीं उबलता: दो किलोमीटर से अधिक जल उस पर
दबाव डालता है, और उस दाब पर जल \qty{100}{\celsius} से बहुत ऊपर भी
द्रव रहता है। काला ``धुआँ'' लोहे के सल्फ़ाइडों और
दूसरे खनिजों के नन्हे कणों का बादल है, जो तब अवक्षेपित होते हैं जब घुली धातुओं से समृद्ध गरम
निकास-जल ठंडे समुद्री जल से मिलता है।
\end{solution}

\begin{solution}{exo:g11:ocean-ridges:6}
(1) लगभग \qty{54}{km} नीचे, लगभग \qty{1324}{\celsius} पर। (2) लगभग
\qty{85}{km} नीचे, लगभग $1400 + 0.45 \times 85 \approx
\qty{1438}{\celsius}$ पर। (3) अधिक गरम प्रावार अधिक गहराई पर पिघलना शुरू करता है, इसलिए
भूपर्पटी तक पहुँचने से पहले ऊपर उठती चट्टान का अधिक लंबा स्तंभ पिघलता है: वह अधिक
बेसाल्ट बनाता है, और अधिक मोटी भूपर्पटी।
\end{solution}

\begin{solution}{exo:g11:ocean-ridges:7}
(1) अपने गलनांक से $1489 - 1345 = \qty{144}{\celsius}$ नीचे।
(2) वह $100 - 54 = \qty{46}{km}$ ऊपर उठती है; अंतर औसतन
प्रति किलोमीटर $144/46 \approx \qty{3.1}{\celsius}$ घटता है। (3) नहीं: ऊपर
जाते हुए चट्टान थोड़ी ठंडी भी होती है, लगभग $0.45 \times 46 \approx
\qty{21}{\celsius}$। गिरता उसका गलनांक है, कहीं अधिक तेज़ी से,
जैसे-जैसे दाब घटता है।
\end{solution}

\begin{solution}{exo:g11:ocean-ridges:8}
अर्ध दरें प्रति वर्ष \qty{7.3}{cm} और \qty{1.65}{cm} हैं, यानी
प्रति मिलियन वर्ष \qty{73}{km} और \qty{16.5}{km}। एक ओर \qty{1000}{km}
तली बनाने के लिए: पूर्वी प्रशांत उभार के लिए $1000/73 \approx 13.7$ मिलियन वर्ष,
मध्य-अटलांटिक कटक के लिए $1000/16.5 \approx 61$ मिलियन वर्ष।
\end{solution}

\begin{solution}{exo:g11:ocean-ridges:9}
(1) \qty{5}{cm} का अर्थ \qty{5e-5}{km} है: $65000 \times 5 \times 10^{-5} =
\qty{3.25}{km^2}$ नई तली प्रति वर्ष। (2) $3.25 \times 7.1 \approx
\qty{23}{km^3}$ भूपर्पटी प्रति वर्ष। (3) $3.6 \times 10^8 / 3.25 \approx 1.1
\times 10^8$ वर्ष, लगभग 110 मिलियन वर्ष: इस गति से कटक
पूरी महासागरीय तली लगभग सौ मिलियन वर्षों में बना सकते हैं, जो पृथ्वी की आयु की तुलना में
थोड़ा समय है। इसीलिए कोई महासागरीय तली बहुत पुरानी नहीं है।
\end{solution}

\begin{solution}{exo:g11:ocean-ridges:10}
(1) भ्रंश-तल क्षैतिज से \ang{60} बनाता है, और धँसाव उस समकोण त्रिभुज की
ऊर्ध्वाधर भुजा है जिसका कर्ण सरकन है:
$1.5/\sin\ang{60} \approx \qty{1.73}{km}$। (2) हर भ्रंश दोनों ओर को
क्षैतिज रूप से $1.5/\tan\ang{60} \approx \qty{0.87}{km}$ दूर करता है; दोनों
भ्रंश मिलकर लगभग \qty{1.73}{km}। (3) नहीं। \qty{25}{km} धँसे खंड को घेरने वाले
भ्रंशों के बीच की दूरी है, खिंचाव की मात्रा नहीं, जो
केवल लगभग \qty{1.7}{km} है; और कटक भ्रंशन के साथ-साथ अक्ष पर नए डाइकों और लावाओं से भी
खुलता रहता है।
\end{solution}

\begin{solution}{exo:g11:ocean-ridges:11}
$d(25) = 6400 - 3200\,e^{-25/62.8} \approx 6400 - 3200 \times 0.672 \approx
\qty{4251}{m}$; $d(100) = 6400 - 3200\,e^{-100/62.8} \approx 6400 - 3200
\times 0.203 \approx \qty{5749}{m}$। 25 और 100 मिलियन वर्षों के बीच
तली लगभग \qty{1500}{m} धँसती है। वर्गमूल नियम 25 मिलियन वर्ष पर
$2500 + 350 \times 5 = \qty{4250}{m}$ देता है: जहाँ दोनों लागू होते हैं, वहाँ दोनों नियम
सहमत हैं।
\end{solution}

\begin{solution}{exo:g11:ocean-ridges:12}
(1) $\alpha\,\Delta T\,L = 3.2 \times 10^{-5} \times 700 \times 100 =
\qty{2.24}{km}$। (2) $s = 3380 \times 2.24/(3380 - 1025) = 7571/2355 \approx
\qty{3.2}{km}$। (3) चरघातांकी नियम कुल $6400 - 2500 =
\qty{3900}{m}$ धँसाव की अनुमति देता है: तापमान-गिरावट और मोटाई के
सरल मानों के साथ, केवल ठंडा होना और सिकुड़ना उसका लगभग
चार-पाँचवाँ भाग समझा देते हैं।
\end{solution}

\begin{solution}{exo:g11:ocean-ridges:13}
(1) $54 - 7 = \qty{47}{km}$ की तुलना में $85 - 7 = \qty{78}{km}$: अधिक गरम
स्तंभ $78/47 \approx 1.66$ गुना लंबा है। (2) $7.1 \times 1.66 \approx
\qty{11.8}{km}$। (3) मापी गई \qty{10.3}{km} सामान्य से अधिक मोटी है, जैसा
पूर्वानुमान था, पर अनुमान से कम: स्तंभ के सबसे गहरे किलोमीटर,
गलन के आरंभ के ठीक ऊपर, उथले किलोमीटरों से कम पिघला पदार्थ देते हैं, इसलिए
प्रति किलोमीटर समान पिघले पदार्थ की धारणा अधिक आकलन करती है।
\end{solution}

\begin{solution}{exo:g11:ocean-ridges:14}
(1) नियम \qty{3178}{m}, \qty{4160}{m}, \qty{4782}{m} और
\qty{5267}{m} देता है; मापी गई माध्यिकाएँ \qty{64}{m},
\qty{142}{m}, \qty{317}{m} और \qty{619}{m} उथली हैं। (2) अंतर आयु के साथ
बढ़ता है। (3) नियम स्वयं तली की गहराई देता है, उसके
अवसादों के लिए संशोधित; मापे गए कोष्ठ अपने अवसाद लिए हैं, जो
पुरानी तली पर अधिक मोटे हैं, और उनमें समुद्री पर्वत और ज्वालामुखीय पठार शामिल हैं, जो
अपने आसपास की तली से उथले हैं।
\end{solution}

\begin{solution}{exo:g11:ocean-ridges:15}
(1) $p = 1025 \times 9.8 \times 2500 \approx \qty{25.1}{MPa}$, और
\qty{0.1}{MPa}: लगभग \qty{25.2}{MPa}। (2) यह क्रांतिक
दाब, \qty{22.064}{MPa}, से ऊपर है: निकास का जल नहीं उबल सकता, चाहे वह कितना भी गरम हो।
(3) दाब \qty{22.064}{MPa} तक इस गहराई पर पहुँचता है:
$(22.064 - 0.1) \times 10^6/(1025 \times 9.8) \approx \qty{2190}{m}$: लगभग
\qty{2200}{m} से गहरा हर निकास उबलने से सुरक्षित है।
\end{solution}

\begin{solution}{pb:g11:ocean-ridges:1}
\textbf{1.} $14.6/2 = \qty{7.3}{cm}$ प्रति वर्ष, यानी प्रति मिलियन वर्ष \qty{73}{km}।

\textbf{2.} कटक-अक्ष की गहराई पर, $d(0) = \qty{2500}{m}$।

\textbf{3.} समुद्री तल से नीचे की ओर: बेसाल्ट लावा (तकिए और चादरी प्रवाह),
चादरी डाइक और \omterm{def:g11:ocean-ridges:gabbro}{गैब्रो}, कुल लगभग \qty{7.1}{km} भूपर्पटी; \omterm{def:g11:seismic-structure:moho}{मोहो} के नीचे
प्रावार का \omterm{def:g11:ocean-ridges:peridotite}{पेरिडोटाइट}, \omterm{def:g11:lithosphere-asthenosphere:lithosphere}{स्थलमंडल} के आधार तक, जो अक्ष पर अभी
बहुत पतला है। अवसाद उस पर बाद में बैठेंगे।

\textbf{4.} जैसे-जैसे प्लेटें एक-दूसरे से दूर होती हैं, \omterm{def:g11:lithosphere-asthenosphere:lithosphere}{दुर्बलतामंडल} कटक के नीचे ऊपर उठता है।
उसका तापमान प्रति किलोमीटर केवल लगभग \qty{0.45}{\celsius} गिरता है,
जबकि उसका गलनांक दाब घटने के साथ कहीं अधिक तेज़ी से गिरता है।
लगभग \qty{54}{km} नीचे दोनों मिलते हैं: चट्टान पिघलना शुरू करती है, और
ऊपर उठते हुए पिघलती रहती है, बिना कोई ऊष्मा जोड़े (\omterm{def:g11:ocean-ridges:partial-melting}{दाब-मुक्ति
गलन})।

\textbf{5.} हर मिलियन वर्ष में कटक अपनी दोनों ओर $2 \times 73 = \qty{146}{km}$
तली बनाता है; कटक के हर किलोमीटर के लिए भूपर्पटी का आयतन
$146 \times 7.1 \times 1 \approx \qty{1040}{km^3}$ है।

\textbf{6.} $73 \times t$: \qty{73}{km}, \qty{292}{km}, \qty{1168}{km} और
\qty{4672}{km}।

\textbf{7.} $2000/73 \approx 27$ मिलियन वर्ष।

\textbf{8.} प्लेटें दसियों मिलियन वर्षों में अपनी चाल और दिशा बदलती हैं,
जैसा हवाई--सम्राट श्रृंखला का मोड़ दिखाता है
(\cref{ch:g11:plate-motions}); आज की दर का 64
मिलियन वर्षों तक बने रहना आवश्यक नहीं है।

\textbf{9.} अक्ष से दूर भूपर्पटी किसी मैग्मा-लेंस के ऊपर नहीं रहती: वह
ठंडी होती है, और नीचे परिसंचरण चलाने और जल को गरम करने के लिए
पर्याप्त ऊष्मा नहीं रहती।

\textbf{10.} पूर्वी प्रशांत उभार के पश्चिम में टुकड़ा प्रशांत
प्लेट पर यात्रा करता है। पश्चिम की ओर प्रशांत प्लेट गहरी खाइयों के अनुदिश समाप्त होती है, जहाँ वह
अपने पड़ोसियों के नीचे डूबती है: क्षेपण क्षेत्र (\cref{ch:g11:subduction-zones})।

\textbf{11.} $d(1) = \qty{2850}{m}$, $d(4) = 2500 + 700 = \qty{3200}{m}$,
$d(16) = 2500 + 1400 = \qty{3900}{m}$, $d(64) = 2500 + 2800 = \qty{5300}{m}$।

\textbf{12.} पहले 16 मिलियन वर्षों में $3900 - 2500 = \qty{1400}{m}$,
और 16 से 64 मिलियन वर्षों के बीच $5300 - 3900 = \qty{1400}{m}$: उतना ही
धँसाव, तीन गुना लंबे समय में। आयु बढ़ने के साथ तली अधिकाधिक धीरे
धँसती है।

\textbf{13.} $350\sqrt{t} = 1500$ से $\sqrt{t} \approx 4.29$ मिलता है, इसलिए
$t \approx 18.4$ मिलियन वर्ष।

\textbf{14.} $6400 - 3200\,e^{-64/62.8} \approx 6400 - 3200 \times 0.361
\approx \qty{5245}{m}$, वर्गमूल नियम से लगभग \qty{55}{m} कम: दोनों
नियम लगभग \qty{1}{\percent} तक सहमत हैं।

\textbf{15.} मापी गई माध्यिका नियम से लगभग \qty{650}{m} उथली है।
नियम अपने अवसादों के भार के लिए संशोधित तली का वर्णन करता है;
मापी गई गहराइयों में अवसाद शामिल हैं, ऐसी तली पर कई दसियों मिलियन वर्षों के
अवसाद, और समुद्री पर्वत और ज्वालामुखीय पठार भी, और ये सब
मापी गई तली को उथला बनाते हैं।

\textbf{16.} $\alpha\,\Delta T = 3.2 \times 10^{-5} \times 700 = 0.0224$:
प्रावार-भाग \qty{2.24}{\percent} अधिक घना है,
$3380 \times 1.0224 \approx \qty{3456}{kg/m^3}$।

\textbf{17.} प्रावार-भाग $100 - 7.1 = \qty{92.9}{km}$ मोटा है। माध्य
घनत्व: $(7.1 \times 2900 + 92.9 \times 3456)/100 \approx (20590 + 321060)/100
\approx \qty{3416}{kg/m^3}$।

\textbf{18.} $3416 > 3380$: पुराना \omterm{def:g11:lithosphere-asthenosphere:lithosphere}{स्थलमंडल} समग्र रूप से अपने नीचे के \omterm{def:g11:lithosphere-asthenosphere:lithosphere}{दुर्बलतामंडल} से लगभग
\qty{36}{kg/m^3}, यानी लगभग \qty{1}{\percent}, अधिक घना है।
हल्की भूपर्पटी अब उसके ठंडे, घने प्रावार-भाग की भरपाई के लिए
पर्याप्त नहीं है।

\textbf{19.} \omterm{def:g11:lithosphere-asthenosphere:lithosphere}{दुर्बलतामंडल} से अधिक घना होने के कारण यह किसी खाई पर वापस
प्रावार में डूब सकता है: क्षेपण, जो अगले अध्याय का विषय है।

\textbf{20.} 64 मिलियन वर्ष पर तली का यह टुकड़ा लगभग
\qty{5300}{m} पर है: लगभग \textbf{5.3~km} नीचे, अपने जन्म-स्थान से
\qty{2.8}{km} अधिक गहरा।
\end{solution}
